%-------------------------------------------------------------------------------
%	SECTION TITLE
%-------------------------------------------------------------------------------
\vspace{1.5mm}
\cvsection{Research Experience}


%-------------------------------------------------------------------------------
%	CONTENT
%-------------------------------------------------------------------------------
\begin{cventries}

    \setlength\leftmargini{4ex}
    \renewcommand{\labelitemi}{---}
    \setlength\leftmarginii{1ex}
    \renewcommand{\labelitemii}{$\circ$}
    \setlength\cventrywidth{2.5cm}
%---------------------------------------------------------

\cventry
{DesCyPhy Lab, Prof. Pierluigi Nuzzo}
{Berkeley, CA}
{Graduate Student Researcher}
{01/2021 -- Present}
{
\vspace{-3.5mm}
\item {\descriptiontitlestyle{Equivalent Neural-Controller Representations}}
\begin{cvitems}
\item Developed a \textbf{provably exact transformation} from discrete-output ReLU neural controllers to soft decision trees, with automatic pruning of redundant branches.
\item Demonstrated verification-time improvements of up to \textbf{20x} on MountainCar, CartPole, and CarRacing; resulted in first-author publications in \textbf{IEEE TAC (2026)} and \textbf{CDC (2023)}.
\end{cvitems}
\vspace{1.8mm}
\item {\descriptiontitlestyle{Compositional Safety for Learning-Enabled CPS}}
\begin{cvitems}
\item Introduced robustness contracts that decouple plant dynamics from neural controllers and enable compositional closed-loop safety verification.
\item Designed SMT- and MILP-based refinement algorithms to tighten abstraction bounds and synthesize realizable assume-guarantee specifications under bounded uncertainty.
\end{cvitems}
\vspace{1.8mm}
\item {\descriptiontitlestyle{Deep RL Controller Distillation for Verifiability}}
\begin{cvitems}
\item Distilled deep-RL policies into compact decision-tree and kernel controllers, improving interpretability and validation while preserving task performance; resulted in an \textbf{Allerton 2022 publication}.
\end{cvitems}
}

\cventry
  {Electronic Design Automation Lab, Prof. Yao-Wen Chang}
  {Taipei, Taiwan}
  {Undergraduate Researcher}
  {09/2018 -- 02/2020}
  {
  \vspace{-3.5mm}
  \item {\descriptiontitlestyle{Routing and Optimization for EDA Contests}}
  \begin{cvitems}
    \item Implemented a multithreaded detailed-routing engine that handled designs with up to \textbf{1M nets} using less than \textbf{64 GB} of memory in 10 hours on most benchmarks.
    \item Built routability-driven global- and system-level FPGA routers; contributed to \textbf{2nd place at ACM ISPD 2019}, a \textbf{Top 5 / Honorable Mention at ICCAD 2019}, and an \textbf{ICCAD 2021 publication}.
  \end{cvitems}
  }

\cventry
  {Institute of Information Science, Academia Sinica}
  {Taipei, Taiwan}
  {Software Researcher}
  {06/2018 -- 09/2018}
  {
    \vspace{-3.5mm}
    \begin{cvitems}
      \item {Implemented a semi-supervised domain-adaptation pipeline for semantic segmentation using unpaired image translation.}
      \item {Improved mean IoU on the Cityscapes benchmark by augmenting segmentation training with CycleGAN-generated domain information.}
    \end{cvitems}
  }

\end{cventries}
